In the context of board games, a common approach for improving model performance is to partition the state space and allocate distinct experts to each region. While this is the central idea behind \emph{Mixture of Experts} (MoE) architectures, the bucketing is often defined manually, requiring domain-specific expertise, and yielding potentially suboptimal partitions.

We investigate and compare several approaches to extending the \emph{Efficiently Updatable Neural Networks} (NNUE) paradigm beyond its conventional formulation, with a particular focus on techniques that have been successfully explored in other areas of machine learning but remain relatively uncommon in NNUE-based chess engines. Mixture-of-Experts architectures are a natural extension for chess evaluation that allows multiple specialized heads to share a common NNUE representation. We also explore sample-gradient-based partitioning, using the gradients of individual training samples to identify groups of positions with similar learning requirements, an approach more commonly studied in large-scale language models and other deep learning settings. Finally, in an effort to make the NNUE architecture more applicable to other fields, we experiment with making its conventional accumulator a learned world model of the game.

The results highlight different trade-offs between the investigated approaches. L1-based clustering combined with a Mixture-of-Experts architecture significantly improves evaluation performance, while introducing only a minor increase in memory usage and inference cost; the lightweight dispatcher accurately reproduces the learned partition. In contrast, sample-gradient-based partitioning is difficult to exploit at inference time, since the gradients depend on the unknown true value of the position, and a lightweight dispatcher struggles to approximate the resulting complex partition. Finally, the world-model approach shows that an NNUE-like architecture can be trained without feedback from a teacher engine or game outcomes, although the unsupervised training objective can reduce evaluation strength. At the same time, learning the world model makes the training process more efficient by shifting a substantial part of the learning effort to the shared representation.
